% Technical development based on IV2025_arun.pdf and the pinned implementation.
\subsection{Dynamic Bicycle Model and Tire Forces}
The full-scale formulation in Sections~III--IV of \cite{muraleedharan2025racing} uses a dynamic bicycle model in a Frenet frame. Each axle represents the combined action of its two tires. Progress $s$, lateral offset $d$, heading error $e_\psi$, body-frame velocities $v_x,v_y$, and yaw rate $r$ describe motion. With path curvature $\kappa(s)$, steering angle $\delta$, mass $m$, yaw inertia $I_z$, and axle distances $l_f,l_r$, its motion equations can be written as
\begin{equation}
\begin{aligned}
\dot s &= \frac{v_x\cos e_\psi-v_y\sin e_\psi}{1-d\kappa},\\
\dot d &= v_x\sin e_\psi+v_y\cos e_\psi,\\
\dot e_\psi &= r-\kappa\dot s,\\
m\dot v_x &= F_{xr}-F_D-F_{yf}\sin\delta\\
 &\quad +F_{xf}\cos\delta+mv_yr,\\
m\dot v_y &= F_{yr}+F_{yf}\cos\delta\\
 &\quad +F_{xf}\sin\delta-mv_xr,\\
I_z\dot r &= l_f(F_{yf}\cos\delta+F_{xf}\sin\delta)\\
 &\quad -l_rF_{yr}.
\end{aligned}\label{eq:fullmodel}
\end{equation}
Here $F_{xi},F_{yi}$ denote longitudinal and lateral axle forces, and $F_D$ denotes aerodynamic drag. The reference article models engine force as a function of throttle, speed, and gear, and models brake force and rolling resistance separately. These powertrain terms belong to its full-scale vehicle model; they are not identified F1TENTH parameters.

The F1TENTH solver uses the reduced state $\mathbf x=[s,d,e_\psi,v_y,r]^\mathsf T$ and scalar input $u=\delta$ \cite{mulakaloori2026repo,muraleedharan2026solver}. The imported solver is pinned at revision \texttt{2c25f4b}. Its forward speed is held constant during each rollout at $\bar v_x=\max(2.5,v_x^{\rm odom})$. It retains the first three equations in Eq.~\ref{eq:fullmodel} and uses
\begin{equation}
\begin{aligned}
\dot v_y &= (F_{yf}\cos\delta+F_{yr})/m-\bar v_xr,\\
\dot r &= (l_fF_{yf}\cos\delta-l_rF_{yr})/I_z.
\end{aligned}\label{eq:reducedmodel}
\end{equation}
The code propagates these states by forward Euler integration, $\mathbf x_{k+1}=\mathbf x_k+\Delta t f(\mathbf x_k,\delta_k)$. Curvature and track widths come from the reference horizon, with desired lateral offset zero. Frenet coordinates require $1-d\kappa\ne0$; the model is intended for motion near the reference path.

Both formulations use a Pacejka-type nonlinear lateral-force law. For the F1TENTH implementation it is
\begin{equation}
\begin{aligned}
F_{yi}&=\mu F_{zi}\sin\!\left[C_i\arctan\!\left(z_i\right)\right],\\
z_i&=B_i\alpha_i-E_i\bigl(B_i\alpha_i-\arctan(B_i\alpha_i)\bigr),\\
\alpha_f&=\arctan\!\frac{v_y+l_fr}{v_{\rm slip}}-\delta,\\
\alpha_r&=\arctan\!\frac{v_y-l_rr}{v_{\rm slip}}.
\end{aligned}\label{eq:tire}
\end{equation}
For the configured zero aerodynamic lift coefficients, $F_{zf}=mg l_r/(l_f+l_r)$ and $F_{zr}=mg l_f/(l_f+l_r)$. Table~\ref{tab:configuration} lists the parameters used in the repository. The reference paper additionally identifies lateral-force offsets and refines its tire model with measured race-car data; those fitted full-scale values are not reused here.

An implementation detail affects Eq.~\ref{eq:tire}: the imported car-steering model uses $v_{\rm slip}=\max(7.0,\bar v_x)$ in m/s. Although the outer F1TENTH header declares a 0.5~m/s constant, it delegates dynamics to the imported model, whose 7.0~m/s floor is active. Thus slip predictions below 7~m/s use a larger denominator than the measured speed. This is a code property of the pinned revision, not a physical tire law, and should be reviewed when tuning low-speed behavior.

\begin{table}[htb]
\centering
\caption{F1TENTH settings from \texttt{config/config.yaml} in \cite{mulakaloori2026repo}. They are configuration values, not parameters measured in the full-scale paper.}
\label{tab:configuration}
\small
\begin{tabular}{lr}
\toprule
Parameter & Value\\
\midrule
Mass $m$ & 3.47 kg\\
Yaw inertia $I_z$ & 0.04712 kg\,m$^2$\\
Front/rear axle distance & 0.15875 / 0.17145 m\\
Wheelbase / body width & 0.3302 / 0.2032 m\\
Friction coefficient $\mu$ & 0.523\\
Front $(B,C,E)$ & $(-6,1.5,-0.1)$\\
Rear $(B,C,E)$ & $(-7,1.5,-0.1)$\\
Steering magnitude limit & 0.4189 rad\\
Steering rate limit & 3.2 rad/s\\
Speed range & 2.5--10.5 m/s\\
Acceleration/deceleration & 7.51 / 8.26 m/s$^2$\\
Stages / samples / basis modes & 100 / 5,000 / 7\\
Rollout step / initial seed & 0.02 s / 70997\\
Warm start & enabled\\
Grid inflation half-width & 4 cells (0.20 m)\\
\bottomrule
\end{tabular}
\end{table}

\subsection{Frequency-Domain Sampling and Input Constraints}
The reference article motivates IDCT sampling because independently sampled time-domain inputs can oscillate excessively \cite{muraleedharan2025racing}. Its algorithm combines steering and longitudinal-force samples before evaluating the vehicle response. The F1TENTH implementation instead evaluates 5,000 scalar steering sequences. The following explicit construction is taken from the imported solver code \cite{muraleedharan2026solver}, supplementing the conceptual description in the article.

Let $N=100$ and $M=7$. Its orthonormal cosine basis is
\begin{equation}
 b_{ki}=\sqrt{2/N}\,c_i\cos\!\left(\frac{\pi i(k+1/2)}{N}\right),
\end{equation}
where $c_0=1/\sqrt{2}$ and $c_i=1$ for $i>0$. For random candidate $j$, coefficients $a_{ij}$ are drawn uniformly in $(-1,1]$. With $h=\dot\delta_{\max}\Delta t$, the increments and bounded commands are
\begin{equation}
\begin{aligned}
\eta_k^{(j)} &=h\sum_{i=0}^{M-1}b_{ki}a_{ij},\\
\delta_k^{(j)} &=\operatorname{clip}_{[-\delta_{\max},\delta_{\max}]}\!\left(
\delta_{k-1}^{(j)}+\operatorname{clip}_{[-h,h]}\eta_k^{(j)}\right).
\end{aligned}\label{eq:idct}
\end{equation}
All candidates start with the current command at $k=0$. The configured rate limit gives $h=0.064$~rad per rollout step. Using low cosine modes restricts rapid variation before clipping. Candidate zero holds the current command; when available, candidate one reuses a shifted previous best sequence under the same limits. This preserves a warm-start candidate without implying global optimality.

\subsection{Tracking, Grip, and Obstacle Costs}
The full-scale article balances tracking error, command variation, terminal error, and tire-grip penalties. In compact notation, its objective has the form
\begin{equation}
\begin{aligned}
J_{\rm IV}={}&\sum_k\left(\mathbf e_k^\mathsf TQ\mathbf e_k+
\Delta\mathbf u_k^\mathsf TR\Delta\mathbf u_k\right)\\
&+\mathbf e_N^\mathsf TQ_f\mathbf e_N
 +\sum_{k,i}q_{g,i}P_{g,i}(k).
\end{aligned}\label{eq:ivcost}
\end{equation}
Here $\mathbf e_k$ denotes tracking error and $\Delta\mathbf u_k$ successive command changes. This notation summarizes Eq.~(11) of \cite{muraleedharan2025racing}. The reference study considers combined force $F_i=\sqrt{F_{xi}^2+F_{yi}^2}$ relative to $\mu_iF_{zi}$ and uses the smooth penalty
\begin{equation}
P_{g,i}=\log\!\left[1+e^{A(F_i/F_{zi}-\mu_{\rm lim})}\right]
-\log\!\left[1+e^{-A\mu_{\rm lim}}\right].
\end{equation}
It penalizes tire-force utilization near the chosen grip threshold. Actuator amplitude and rate bounds restrict sampled commands, whereas a finite grip penalty does not guarantee constraint satisfaction.

The actual F1TENTH cost differs from Eq.~\ref{eq:ivcost}. Define the error penalty $E(z;q,c,l)=\min(qz^2,c)+l|z|$. With $w_k=(1-0.3k/(N-1))/N$, the implementation evaluates
\begin{equation}
\begin{aligned}
J_{\rm F1}={}&\sum_{k=0}^{N-1}\left(w_kL_k+
\frac{20000}{N}\mathbb I_{\rm occ}(\mathbf x_k)\right)\\
&+\frac{E_{d,T}(d_{N-1})+E_{\psi,T}(e_{\psi,N-1})}{N},\\
L_k={}&E_d(d_k)+E_\psi(e_{\psi,k})+0.05\delta_k^2-0.5s_k\\
&+5(\delta_k-\delta_{k-1})^2
 +P_{f,k}+P_{r,k}+P_{b,k}+P_{\alpha,k}.
\end{aligned}\label{eq:f1cost}
\end{equation}
The command-change term is zero at $k=0$. The $(q,c,l)$ triples are $(400,600,50)$ for lateral error and $(3500,250,500)$ for heading error; terminal triples are $(200,400,20)$ and $(1500,200,100)$. The negative progress term rewards forward motion. The force penalties $P_f,P_r$ use lateral-force utilization $z_i=|F_{yi}|/(\mu F_{zi})$ and the barrier $2\{\log[1+\exp(20(z_i-1))]-\log[1+\exp(-20)]\}$. In particular, the configured pair \texttt{[20,2]} maps to sharpness 20 and multiplier 2 in the imported structure.

For each track side, $P_b$ adds a linear penalty of gain 1500 for exceeding the boundary after subtracting half the car width and an extra 0.08~m margin, capped at 1000 per side. The slip penalty adds $\min(5000,1000\max(0,|\alpha_i|-0.15))$ per axle. These terms express preferences through finite costs.

The occupancy indicator in Eq.~\ref{eq:f1cost} is obtained by projecting a predicted Frenet state into the local LiDAR grid, using curvature at the current vehicle position. A hit in an inflated occupied cell adds $20000/N=200$ per stage, independently of $w_k$. LiDAR returns are inflated by a square extending four cells in each direction. The lookup treats positions outside the local grid as unoccupied, so grid extent and the constant-curvature projection limit what this obstacle cost observes. The objective minimizes cost over the sampled candidates; it supplies neither a continuous-optimization optimum nor a formal collision-avoidance guarantee.

\subsection{Receding-Horizon Execution on the GPU}
The implementation follows the sampling, prediction, evaluation, and selection sequence of Algorithm~1 in \cite{muraleedharan2025racing}, specialized to steering:
\begin{enumerate}
\item Convert odometry to the five-state Frenet model and construct the reference horizon. Update the local grid from available LiDAR scans.
\item Generate bounded IDCT steering sequences, including the hold and warm-start candidates.
\item Propagate each candidate through Eqs.~\ref{eq:reducedmodel}--\ref{eq:tire} and accumulate Eq.~\ref{eq:f1cost}. CUDA executes candidate rollouts in parallel, with sequential state propagation within each candidate.
\item Copy candidate costs to the host and select the smallest finite cost. Copy the winning trajectory and retain it for the next warm start.
\item Publish the winning sequence's element at index 1, since index 0 is the current command. Compute the speed command separately and repeat at the next odometry callback.
\end{enumerate}
The node publishes zero speed and steering if the solver returns an error. The initial random seed is 70997 and is incremented by the solve-call index. Consequently, reproducing the sampled sequence also depends on the sequence of callbacks. Unlike a wholly GPU-resident minimum-cost reduction, this pinned solver selects the winner on the CPU after copying all candidate costs. This distinction matters when interpreting timing relative to the GPU-focused full-scale study.

The separate speed command uses Eq.~\ref{eq:speed}, a forward-grid-distance heuristic, and clipping to the configured speed range. Its slew limiter advances by $a_{\max}\Delta t$ or $b_{\max}\Delta t$ per callback using the configured rollout step, rather than measured callback spacing. Therefore the tabulated acceleration settings are nominal increments in that implementation; their realized per-second effect depends on the callback rate. The speed floor is not a stopping policy near obstacles. These details should be included when reproducing or modifying the controller.
